\section{Le système mécanographique~: le couplage stockage-calcul (1890--1940)}\label{sec:mechanography}
% =============================================================================

La calculatrice mécanique accélérait un geste~: un opérateur entrait des nombres, la machine produisait un résultat, rien ne persistait. Le résultat apparaissait sur un cadran ou sur une bande de papier que l'opérateur recopiait dans un registre, et la machine n'en gardait aucune trace. Chaque calcul était un acte isolé, sans mémoire, sans lien avec le calcul précédent ni avec le suivant. C'est la raison pour laquelle cent firmes pouvaient coexister sur ce marché~: une machine Burroughs et une machine Brunsviga faisaient la même chose, et aucune des deux ne créait de dépendance au-delà de l'habitude de l'opérateur.

Ce que Herman Hollerith construisit à la fin des années~1880 était d'une nature différente. La carte perforée~--- un rectangle de carton portant des trous disposés selon un code~--- était un support de données séparé de l'opérateur et de la machine, réutilisable, triable, stockable. Elle pouvait être perforée par une perforatrice, triée par une trieuse, tabulée par une tabulatrice, reproduite par une reproductrice~: le même jeu de cartes passait d'une machine à l'autre dans un flux intégré. Hollerith réalise ainsi, pour la première fois dans un système opérationnel, le couplage entre un stockage sur support \emph{machine-lisible} et un traitement \emph{en volume} de l'information qui y est inscrite. La précision n'est pas verbale. Le métier Jacquard (1804) couplait certes une carte perforée à un dispositif mécanique, mais la chaîne de cartes y constituait \emph{un programme unique} exécuté de façon répétée sur un substrat continu (le tissu)~--- non un ensemble de \emph{données distinctes} traitées par un opérateur agrégateur. L'\emph{Analytical Engine} de Babbage, conçue dès les années~1830, séparait le \emph{store} et le \emph{mill} dans une architecture analogue à celle de Hollerith, mais n'avait jamais été réalisée. Les tables de Prony, enfin, couplaient un calcul humain divisé à un stockage papier réutilisable, mais le support n'était pas lisible par une machine. Aucun de ces dispositifs ne satisfait la combinaison \emph{stockage machine-lisible / traitement en volume de données distinctes / système opérationnel} qui définit le geste hollerithien~--- et c'est cette combinaison, non l'un quelconque de ses composants, qui restructure ce qui suit. Ce n'était plus un outil qui accélérait un geste~; c'était un système qui restructurait un processus.

\textcite{cortada_before_1993} l'a formulé avec netteté~: \og \emph{unlike typewriters or calculating machines, which were single units of equipment that handled small amounts of information in a restricted fashion, tabulating gear operated on thousands, even millions of pieces of data}\fg{}. La distinction entre l'outil isolé et le système intégré n'est pas une question de degré mais de nature. Elle conditionne tout ce qui suit~--- le modèle économique, la forme du monopole, le rapport au travail, et la trajectoire qui mène à l'ordinateur.


\subsection{La fiche comme infrastructure~: le stockage avant les machines}\label{subsec:fiche_infrastructure}

Avant que la tabulatrice ne couple le stockage au calcul, le stockage existait déjà comme dispositif technico-organisationnel autonome. Dans la seconde moitié du XIX\textsuperscript{e}~siècle, les bibliothèques, les administrations, les bureaux de banque et les services médicaux développent un appareil~--- la fiche, le classeur, le meuble à tiroirs~--- qui constitue, indépendamment de toute mécanisation, une infrastructure de l'information. \textcite{gardey_ecrire_2008} a montré que cet appareil ne se réduit ni à ses objets matériels, ni à ses modes d'usage individuels~: il est d'emblée un \emph{système}, dont les pièces ne prennent leur sens que dans l'ensemble qui les ordonne. \og Une, la fiche n'est rien~; elle ne prend son sens que conçue dans cet ensemble (non fini) dont elle peut être extraite\fg{}\footnote{\textcite[ch.~4]{gardey_ecrire_2008}, qui en fait la \og lecture forte\fg{} du dispositif. La formule restitue ce que les manuels d'organisation du travail de bureau de l'époque~--- \textcite{leffingwell_office_1926} aux États-Unis, \textcite{bolle_utilisation_1929} en France~--- donnent comme évidence pratique~: une fiche utilisée seule est un fragment, et l'investissement n'a de rendement que s'il porte sur le système complet, c'est-à-dire l'index, les meubles, les protocoles d'inscription et les opératrices formées.}.

\dimmark{D8}{*}%
Cette transformation, datée par Gardey de la période~1890--1920, reconfigure ce que \textcite{febvre_outillage_1938} nommera dans les années~1930 l'\og outillage mental\fg{} d'une époque~: non plus seulement les concepts et les mots disponibles à la pensée, mais les supports, les instruments et les gestes qui rendent certaines opérations cognitives possibles\footnote{La filiation est explicitement revendiquée par \textcite[introduction]{gardey_ecrire_2008}, qui propose son ouvrage comme \og prolongation matérielle\fg{} de la notion d'outillage mental introduite par Febvre dans les premiers \emph{Annales}. La présente section retient cette extension matérielle de l'outillage mental comme cadre opératoire pour l'analyse de la fonction STORE~: ce ne sont pas seulement des objets que l'époque produit, ce sont des dispositifs distribués qui rendent disponibles des opérations~--- classer, retrouver, croiser, comparer~--- que les supports manuscrits antérieurs ne supportaient qu'au prix d'un travail considérable.}. La fiche, dans cette perspective, n'est pas un perfectionnement marginal du registre manuscrit~: elle inaugure une économie matérielle de l'écrit dans laquelle la mobilité, la malléabilité et la combinatoire prennent le pas sur la sécurité et la conservation. La confiance que les pratiques anciennes plaçaient dans les humains qualifiés~--- le commis qui retrouve un dossier, le bibliothécaire qui connaît son fonds~--- est progressivement déléguée aux artefacts et aux protocoles qui en organisent le repérage.

L'autre vertu théorique de la fiche, dont la sociologie des sciences a tiré des conséquences, est sa plasticité interprétative. \textcite{star_griesemer_1989} ont qualifié d'\og objets-frontières\fg{} (\emph{boundary objects}) les artefacts qui circulent entre des communautés de pratiques distinctes en gardant assez de stabilité pour permettre la coordination, mais en restant assez plastiques pour servir des usages différents. La fiche médicale, la fiche bibliographique, la fiche actuarielle, la fiche de stocks~: ces objets matériellement comparables soutiennent des opérations cognitives radicalement différentes selon le système qui les ordonne. Ce qui les rend \emph{infrastructure} n'est pas leur identité matérielle mais leur capacité à devenir le substrat partagé de coordinations hétérogènes\footnote{La théorie des objets-frontières de Star et Griesemer est postérieure aux faits que cette section décrit~; elle est ici employée rétrospectivement comme grille de lecture, dans la logique d'anachronisme contrôlé exposée au préambule du chapitre.}.

\dimmark{D7}{+}%
Cette infrastructure-fiche fonctionne déjà à grande échelle bien avant que la mécanisation ne s'en saisisse. Au tournant du siècle, le bibliographe belge Paul~Otlet et son associé Henri~La~Fontaine entreprennent de constituer, à Bruxelles, un \emph{Répertoire bibliographique universel} dont l'ambition est de fournir un index unifié de l'ensemble du savoir publié~; ils accumuleront, jusqu'à la Première Guerre mondiale, plusieurs millions de fiches selon une classification décimale qu'ils ont eux-mêmes affinée\footnote{Sur le projet d'Otlet et de La~Fontaine et leur \emph{Mundaneum} bruxellois, la référence biographique et archivistique canonique reste \textcite{rayward_universe_1975}.}. Otlet, dans les textes de cette période repris dans son \emph{Traité de documentation}~\parencite{otlet_traite_1934}, emploie déjà l'expression de \og société de l'information\fg{} pour désigner le régime que ce dispositif rend possible\footnote{Sur cette occurrence pré-1914 du syntagme et sur le projet plus large d'une économie politique de la documentation, on renverra à l'enquête d'Anne~Rasmussen reprise par \textcite[ch.~7]{gardey_ecrire_2008}~: l'expression circule dans les milieux internationaux de la documentation scientifique et technique avant d'être réinventée, sous une autre généalogie, par \textcite{machlup_production_1962}, qui en méconnait la profondeur historique.}. Que le concept précède de plus d'un demi-siècle l'industrie qui s'en réclamera ne tient pas à une anticipation visionnaire~: il tient à ce que l'infrastructure-fiche, autonome, suffisait déjà à poser comme objet le traitement systématique de l'information.

C'est sur cette infrastructure préexistante que vient se greffer la mécanographie. La carte perforée d'Hollerith n'invente ni la fiche, ni le système qui l'ordonne, ni l'opératrice qui la manipule~: elle couple ces éléments à un dispositif de calcul électromécanique. Ce que la sous-section suivante examine n'est donc pas l'irruption d'une fonction nouvelle~--- la fonction STORE existe avant elle~--- mais la transformation que produit son couplage à la fonction COMPUTE. La singularité du moment Hollerith-IBM tient précisément à ce couplage~: le stockage, déjà infrastructure, devient \emph{lisible par une machine}, et le calcul, déjà industriel, devient capable de \emph{traiter le stocké en volume}.

\emergence{STORE~--- La fonction de stockage de l'information n'apparaît pas avec la mécanographie~: elle est constituée comme infrastructure technico-organisationnelle dès la seconde moitié du XIX\textsuperscript{e}~siècle, autour de la fiche, du classeur et des protocoles administratifs et bibliographiques qui les ordonnent. Ce que la mécanographie introduit n'est pas la fonction mais son couplage avec la fonction de calcul. Ce déplacement~--- du stockage comme infrastructure manuelle à son intégration dans une chaîne automatisée~--- est l'objet propre de la présente section. \questionch{2}{D7/D8~--- La distinction entre l'infrastructure-fiche pré-mécanographique et son couplage avec le calcul électromécanique permet-elle de définir un seuil structurel~? Si oui, ce seuil est-il l'analogue, pour STORE, de ce qu'est l'unification de marché par le télégraphe pour TRANSMIT~?}}


\subsection{Hollerith et le recensement de~1890}\label{subsec:hollerith}

Herman Hollerith (1860--1929) était employé au \emph{Bureau of the Census} quand, en~1881, John Shaw Billings, directeur de la division des statistiques démographiques, lui fit la remarque qui allait orienter sa carrière~: il devait bien exister un moyen mécanique de tabuler les données du recensement, quelque chose sur le principe du métier Jacquard \parencite{cortada_before_1993}. Le recensement de~1880, portant sur cinquante millions d'habitants, avait pris huit ans à dépouiller~; celui de~1890, avec une population estimée à plus de soixante millions et un questionnaire considérablement élargi, menaçait de ne pas être achevé avant le suivant.

Hollerith construisit, entre~1884 et~1889, un système complet~: une perforatrice manuelle pour inscrire les données sur des cartes, un trieur électrique pour regrouper les cartes par catégories, et une tabulatrice pour compter les trous. La tabulatrice de~1890 ne faisait qu'une chose~--- compter~--- mais elle le faisait vite~: 415~cartes par minute, soit 80\,000 à 90\,000 par jour \parencite{cortada_before_1993}. En~1890, ses machines tabulèrent les 62\,622\,250~habitants des États-Unis et permirent de publier les premiers résultats en six semaines, contre des années par les méthodes antérieures. Le \emph{Bureau of the Census} estima l'économie à cinq millions de dollars par rapport au dépouillement manuel. Pour le recensement de~1900, le Bureau utilisa 311~tabulatrices, 20~trieuses automatiques et 1\,021~perforatrices, pour un coût total de 428\,239~dollars incluant les machines, le service et les cartes \parencite{cortada_before_1993}.

Le succès du recensement américain ouvrit les portes internationales~: l'Autriche utilisa les machines Hollerith pour son recensement de~1891, le Canada et la Norvège suivirent, puis la Russie en~1897 pour ses 129~millions de citoyens \parencite{cortada_before_1993}. Mais le recensement était un événement décennal~; ce dont Hollerith avait besoin pour pérenniser son entreprise, c'était d'une base de clientèle continue. Il la trouva dans les chemins de fer~--- le \emph{New York Central Railroad} signa en~1895 un contrat pour traiter près de quatre millions de lettres de voiture par an~--- puis dans les assureurs, les sidérurgistes et les grands distributeurs. Avant~1914, les applications couvraient la paie, la gestion des stocks, la comptabilité analytique, le suivi des coûts de production~--- toutes les fonctions que les calculatrices mécaniques ne pouvaient pas traiter en volume.


\subsection{Le modèle économique~: location, cartes, brevets}\label{subsec:ibm_model}

\dimmark{D2}{*}%
Le modèle économique d'Hollerith, puis de C-T-R et d'IBM, était structurellement différent de celui des fabricants de calculatrices. Les machines n'étaient pas vendues mais louées~: une trieuse coûtait environ 10~dollars par mois, une tabulatrice 40~dollars \parencite{cortada_before_1993}. L'amortissement intervenait en vingt mois pour les tabulatrices, trente pour les trieuses. Mais la vraie source de profit n'était pas la location des machines~--- c'étaient les cartes. Le contrat de location stipulait que le client devait acheter ses cartes exclusivement auprès d'Hollerith (puis d'IBM). Les cartes coûtaient entre 85~cents et un dollar le millier~; leur production revenait à environ 30~cents, soit une marge de 60 à 80\,\% \parencite{cortada_before_1993}. En~1929, les ventes de cartes représentaient un tiers des revenus d'IBM~; pendant la dépression, cette proportion monta à 39\,\%, parce que les clients qui gardaient leurs machines devaient continuer à acheter des cartes. Au pire de la crise, IBM vendait cent millions de cartes par an à 1,40~dollar le millier. En~1938, les cartes rapportaient 5~millions de dollars sur 34,7~millions de revenus totaux et étaient le produit le plus rentable de la gamme \parencite{cortada_before_1993}.

Le mécanisme de monopole qui en résultait procédait du verrouillage du consommable~: la machine était un véhicule pour la vente de cartes, et le client qui avait perforé ses données sur des cartes IBM ne pouvait pas les faire lire par des machines Powers\footnote{Les cartes IBM à~80~colonnes (1928) utilisaient des perforations rectangulaires, délibérément différentes des perforations rondes des cartes Hollerith antérieures et des cartes Powers. \textcite{heide_punched_2009} note que les perforations rectangulaires furent choisies non seulement pour des raisons techniques (lecture électrique plus fiable, carton plus solide) mais aussi parce qu'elles étaient brevetables.}. \textcite{austrian_hollerith_1982} montre que Hollerith avait imposé cette clause d'exclusivité dès les premiers contrats de location au \emph{Bureau of the Census}, et que le modèle fut maintenu systématiquement par Watson après la création de C-T-R. À cela s'ajoutait un portefeuille de brevets systématiquement élargi par des rachats~: IBM acquit les brevets de Peirce en~1922 pour un montant équivalant à 25\,\% de ses bénéfices nets annuels, et racheta en~1930 l'ensemble des brevets de l'Autrichien Gustav Tauschek~--- non pour utiliser sa technologie, mais pour empêcher un concurrent de contrôler la multiplication par cartes perforées \parencite{heide_punched_2009}.

Les chiffres du \emph{Department of Justice} américain, à la fin de~1935, quantifient le résultat~: IBM contrôlait 85,7\,\% des tabulatrices installées aux États-Unis, 86,1\,\% des trieuses, 81,6\,\% des perforatrices \parencite{cortada_before_1993}. L'antitrust américain attaqua IBM et Remington Rand en~1933, et la Cour suprême confirma en~1936 l'interdiction d'obliger les clients à acheter les cartes auprès du loueur. Mais le monopole survécut au jugement~: la dépendance technique créée par les formats de cartes propriétaires et la supériorité de la force de vente IBM suffirent à maintenir la position dominante.

\dimmark{D2}{*}%
\textcite{heide_punched_2009} éclaire un aspect que les données américaines seules ne font pas voir. La licence que Watson accorda à Powers en~1914, donnant accès aux brevets Hollerith contre une redevance de 20\,\% du chiffre d'affaires sur les tabulatrices et trieuses, était un instrument de contrôle autant qu'un geste de prudence antitrust. En Allemagne, un tribunal avait imposé une redevance de 5\,\% seulement pour une licence équivalente \parencite{heide_punched_2009}. Le différentiel~--- 20\,\% aux États-Unis, 5\,\% en Allemagne~--- illustre une différence institutionnelle que Heide formule explicitement~: le système américain de brevets renforçait le premier entrant en lui permettant de fixer les conditions de la licence, tandis que le système allemand, avec ses licences imposées par voie judiciaire, facilitait la concurrence.

\textcite{mounier-kuhn_punched_2015}, dans sa recension de Heide, propose une périodisation en quatre cycles de clôture (\emph{closure}) du système des cartes perforées, chacun fondé sur un brevet-clé renouvelé~: le brevet Hollerith initial (lecture électrique), le brevet sur le groupe automatique (1914, accordé en~1931 seulement), le brevet sur la carte à~80~colonnes (1928), et les brevets de multiplication. À chaque cycle, le renouvellement du brevet-clé permettait de prolonger le verrouillage du système et de bloquer ou de retarder les challengers.

\emergence{D2~--- Le monopole IBM procède d'un mécanisme différent de celui du télégraphe~: pas les rendements de réseau mais le triptyque location + cartes exclusives + brevets. Deux monopoles, deux mécanismes, un même résultat~--- la concentration du marché dans les mains d'un opérateur dominant. \questionch{2}{D2~--- Le monopole par verrouillage du consommable (IBM/cartes) est-il structurellement plus stable que le monopole par rendements de réseau (télégraphe)~?}}

\dimmark{D3}{+}%
Les revenus d'IBM passèrent de 2,5~millions de dollars en~1919 à 20,3~millions en~1931, fléchirent à 17,6~millions en~1933, puis remontèrent à 45,3~millions en~1940 grâce aux contrats du New Deal \parencite{cortada_before_1993}. Le profit net passa de 1,4~million en~1922 à 9,1~millions en~1939, et IBM paya des dividendes sans interruption pendant toute la dépression. En~1928, IBM n'était que le quatrième fournisseur de l'industrie des machines de bureau par le chiffre d'affaires (19,7~millions, derrière Remington Rand à 59,6, NCR à 49, et Burroughs à 32,1), mais le premier par la marge~: 26,9\,\% de profit sur le chiffre d'affaires, contre 25,9\,\% pour Burroughs, 15,9\,\% pour NCR et 10,1\,\% pour Remington Rand \parencite{cortada_before_1993}. En~1939, les positions s'étaient inversées~: IBM était premier en profit (9,1~millions), devant NCR (3,1) et Burroughs (2,9), alors que Remington Rand, avec un chiffre d'affaires encore supérieur (43,4~millions), ne dégageait que 1,6~million de bénéfice. C'est la marge qui distinguait IBM, pas la taille~--- et cette marge venait des cartes.

Pour autant, la taille relative d'IBM dans l'économie américaine restait modeste. En~1930, IBM représentait 12\,\% de la valeur de l'industrie des machines de bureau (20,3~millions sur 165,3) et moins de trois centièmes de pour cent du produit national brut. L'industrie entière~--- calculatrices, tabulatrices, machines à écrire, caisses enregistreuses~--- pesait 217,8~millions de dollars en~1929, soit un peu plus de deux dixièmes de pour cent du PNB \parencite{cortada_before_1993}. Watson lui-même reconnaissait en~1929 que son équipement n'était installé que dans un cinquième des organisations qui pourraient l'utiliser~--- le marché potentiel n'était servi qu'à la marge. La raison en était le seuil de rentabilité~: un client devait effectuer au moins 1\,500~opérations par jour pour justifier l'investissement \parencite{cortada_before_1993}, ce qui excluait de fait les petites organisations.


\subsection{Une trajectoire alternative~: les classi-compteurs de Lucien March (1901--1940)}\label{subsec:march_classicompteurs}

En~1896, lorsque Lucien March (1859--1933) prend la direction de la Statistique générale de France, le bureau parisien dépouille le recensement quinquennal au moyen des machines Hollerith\footnote{Polytechnicien de formation, Lucien March dirige la Statistique générale de France de~1896 à~1920. Sur sa carrière administrative et son rôle dans la modernisation statistique française, voir \textcite{cinquanteans_insee_1996} et, pour le contexte plus large de l'épistémologie statistique française, \textcite{desrosieres_politique_1993}.}~; en~1901, il ne le fait plus. Le bureau s'en décharge à l'occasion d'un dispositif que March a conçu lui-même et que la Statistique utilisera, sans interruption majeure, pour tous les recensements jusqu'à celui de~1936. March avait perçu, dès l'expérience du recensement de~1896, deux limites au système américain~: son coût~--- les machines Hollerith étaient louées pour des sommes importantes, et leur usage continu n'était pas justifié par le rythme décennal du recensement~--- et la séparation qu'il imposait entre la lecture des bulletins et le comptage statistique, en obligeant à une perforation préalable comme étape distincte du dépouillement.

Le classi-compteur, breveté par March à la fin du XIX\textsuperscript{e}~siècle et progressivement perfectionné, est un clavier compté de soixante touches, chacune correspondant à une donnée portée sur le formulaire de recensement~--- sexe, date et lieu de naissance, profession, situation matrimoniale, et ainsi de suite. L'opératrice~--- les sources administratives parlent des \og dames\fg{}~--- enfonce successivement les touches en lisant chaque bulletin, et les compteurs internes additionnent par catégorie. Une fois la série traitée, l'abaissement d'un châssis imprime sur une bande les chiffres enregistrés, le compteur est remis à zéro et un nouveau lot peut être dépouillé. Le rythme de fonctionnement est de huit mille bulletins par jour \parencite[ch.~7]{gardey_ecrire_2008}.

\dimmark{D10}{*}%
L'innovation analytique du dispositif tient à un geste~: March élimine la carte perforée comme support intermédiaire. Là où le système Hollerith requiert une perforation préalable de la donnée sur un support indépendant, qui doit ensuite être trié et tabulé en flux successifs, le classi-compteur fait coïncider la lecture du bulletin avec son comptage. Le bulletin papier reste comme document source, mais l'objet-fiche-perforée disparaît du processus~--- ce qui supprime à la fois le poste budgétaire le plus coûteux du système américain (la location des perforatrices, l'achat des cartes) et l'étape de travail la plus exposée à l'erreur de transcription. L'inconvénient inverse~--- l'absence d'un support intermédiaire matériel qui permette une vérification à distance et a~posteriori~--- pose en contrepartie des problèmes en matière de détection d'erreurs, mais sans remettre en cause le choix du dispositif pour l'usage ciblé~: le dépouillement statistique d'État.

\dimmark{D2}{*}%
La Statistique générale de France utilise les classi-compteurs pour les recensements de~1901, 1906, 1911, 1921, 1926, 1931 et~1936, avec des perfectionnements successifs portant sur la mécanique du clavier et l'extension du jeu de catégories codables. L'avantage économique est documenté~: moins cher que les locations Hollerith, le système permet par ailleurs aux opératrices de manipuler directement les bulletins sans étape de perforation séparée. Au-delà du cas français, le système est adopté par les administrations statistiques d'Italie, du Portugal et du Brésil~\parencite[ch.~7]{gardey_ecrire_2008}~; \textcite{peaucelle_mecanographie_2004}, dans une étude économique rétrospective, juge que le choix de la mécanographie classi-compteurs en France relève d'une rationalité économique défendable au regard des coûts d'opportunité de l'époque, et non d'un retard ou d'une singularité culturelle.

Ce cas mérite d'être inscrit avec soin dans la lecture comparative que cette section conduit. Il ne s'agit pas d'une variante secondaire du système Hollerith~: c'est un dispositif structurellement différent~--- le clavier-compteur sans support intermédiaire~--- qui a fonctionné quarante ans dans plusieurs pays, et qui n'a pas été supplanté par les cartes perforées dans son espace d'usage. Ce que ce cas révèle, c'est que la trajectoire Hollerith-IBM, qui domine l'historiographie de la mécanographie, n'était ni la seule disponible ni la seule économiquement viable au tournant du siècle. La concentration de l'historiographie sur Hollerith tient en partie à la croissance ultérieure du conglomérat américain, mais cette croissance ne peut être rétrojetée comme nécessité technique~: à l'époque où le système IBM se déployait aux États-Unis et conquérait progressivement les administrations statistiques européennes, en France, le classi-compteur dépouillait les recensements à un coût inférieur, et continuerait à le faire pendant quatre décennies. La non-émergence d'IBM en France pour le dépouillement statistique d'État n'est donc pas un retard~; c'est une trajectoire alternative documentée. \textcite{mounier_kuhn_informatique_1999} insiste sur ce point dans son enquête sur l'informatique française d'avant le \emph{Plan calcul}~: les institutions statistiques françaises ne se sont rangées au modèle des cartes perforées qu'au sortir de la Seconde Guerre mondiale, et seulement parce que le volume et la complexité du dépouillement~--- introduction de variables nouvelles, croisements multivariés~--- débordaient désormais les capacités du clavier-compteur.

\emergence{D2/D10~--- Le cas March établit empiriquement qu'à demande équivalente (le dépouillement décennal d'un recensement national), plusieurs dispositifs technico-organisationnels coexistent durablement et occupent la même fonction par des architectures matérielles distinctes. La domination ultérieure d'IBM ne relève donc pas d'une supériorité technique intrinsèque mais d'un choix d'écosystème institutionnel et économique. \questionch{2}{D2/D10~--- Quel mécanisme institutionnel a permis à la trajectoire alternative française de se maintenir quarante ans, et qu'est-ce qui a fait qu'elle ne s'est pas exportée au-delà de quelques pays latins~? Le facteur est-il l'absence d'industrie de fabrication concurrente, l'absence de stratégie d'exportation comparable à celle de Watson, ou la spécificité du dépouillement administratif comme niche fonctionnelle relativement à d'autres usages~?}}


\subsection{La tabulatrice dans les organisations~: le Nautical Almanac et le chemin de fer}\label{subsec:zoom_tabulating}

Le Census de 1890 avait démontré que la tabulatrice pouvait traiter des millions d'enregistrements en un temps que le calcul humain ne pouvait pas atteindre. Mais le recensement est un événement décennal~; ce qui allait révéler les effets de la mécanisation sur l'organisation du travail, c'est l'usage quotidien de ces machines dans des institutions qui calculent en continu~--- les observatoires, les compagnies d'assurance, les administrations ferroviaires. C'est dans ces usages que la tabulatrice cesse d'être un exploit technique et devient un fait organisationnel, dont les conséquences sur le travail, le coût et la supervision n'avaient pas été anticipées par ses promoteurs.

\paragraph{Le Nautical Almanac~: la mécanisation et ses coûts}\mbox{}\\
Le cas le mieux documenté est celui du \emph{Nautical Almanac} britannique, dont les archives administratives ont été exploitées par \textcite{daston_calculation_2018} à partir du fonds RGO~16 de la Cambridge University Library. Le \emph{Nautical Almanac}, publié depuis 1767 sous la direction de l'\emph{Astronomer Royal}, fournissait les éphémérides nécessaires à la navigation de la flotte marchande et militaire britannique. Sa production exigeait des dizaines de milliers de calculs répétitifs~--- positions lunaires, prédictions de marées, tables de réfraction~--- dont l'exécution avait été organisée, depuis Maskelyne, en un réseau de calculateurs à domicile dispersés dans toute l'Angleterre \parencite{croarken_human_2009}. Chaque calculateur recevait des \og \emph{precepts}\fg{}~--- des algorithmes décomposés en étapes sur des formulaires pré-imprimés~--- et rendait un mois complet de tables, vérifié par un \emph{anti-computer} et un \emph{comparer}. Le système relevait du \emph{putting-out}, comme le tissage à domicile~: du travail qualifié, décentralisé, faiblement supervisé.

En décembre~1928, l'Amirauté autorisa l'achat d'une machine Burroughs et la location d'une machine Hollerith pour une durée de six mois. Le surintendant Comrie justifia la dépense par la promesse d'une réduction de la fatigue mentale et d'une augmentation de la précision. Les archives permettent de reconstituer le bilan. La location de la Hollerith coûtait environ 264~livres~; les dix mille cartes perforées nécessaires, 100~livres~; les salaires de quatre opératrices pendant six mois, puis de deux opératrices pendant six mois supplémentaires, 234~livres~; l'électricité, 9~livres. Le total atteignait 607~livres pour un calcul que les méthodes anciennes accomplissaient pour 500~livres par an. La mécanisation coûtait plus cher, pas moins. Et elle exigeait davantage de travail, pas moins~: au lieu de dix mille sommes tirées de sept tables différentes, il fallait désormais perforer douze millions de chiffres sur trois cent mille cartes pour les faire passer dans la machine.

\dimmark{D4}{*}%
Ce qui avait changé, ce n'était pas le volume de travail mais sa nature et sa répartition. Les anciens calculateurs~--- des hommes qualifiés, souvent retraités du service ou enseignants complétant leurs revenus, travaillant chez eux à leur rythme~--- furent remplacés par des opératrices recrutées sur concours dans les catégories les moins bien payées de la fonction publique. L'Amirauté exigeait des candidates qu'elles eussent passé un examen compétitif portant sur \og \emph{English, Arithmetic, General Knowledge and Mathematics}\fg{}~; la réglementation de la \emph{Civil Service} interdisait l'emploi de femmes mariées. Le surintendant Comrie, dans une lettre de 1931, estimait que des femmes pourraient occuper les postes d'assistante aussi bien que des hommes, mais recommandait que les fonctions de surintendant et de surintendant adjoint \og \emph{be reserved for men, especially in view of the fact that the greater part of the calculation is now performed by mechanical means}\fg{}~--- logique remarquable où la mécanisation du calcul, loin de rendre le genre du calculateur indifférent, justifiait au contraire la réservation de l'autorité sur les machines aux hommes qui ne les opéraient pas.

La centralisation devint obligatoire~: les machines étaient coûteuses et ne pouvaient pas quitter le bureau. Le travail à domicile disparut. Et le travail de supervision explosa. Là où l'ancien système consistait à dire à un calculateur~--- un certain W.~F.~Doaken, \emph{M.A.}~--- \og \emph{Do the Moon Transit}\fg{}, et à recevoir cinq mois plus tard la copie prête pour l'imprimeur, le nouveau système distribuait le travail entre six ou sept personnes, auxquelles le surintendant devait fournir \og \emph{perhaps 100 to 120 different sets of instructions}\fg{}. \dimmark{D7}{+}%
Comrie estimait que la préparation du travail pour les machines et leurs opératrices représentait désormais vingt à trente pour cent du calcul total~--- un cinquième à un tiers du temps du surintendant consacré non pas à calculer, ni à superviser au sens classique, mais à \emph{concevoir la séquence} dans laquelle humains et machines devaient s'articuler. C'est ce que \textcite{daston_calculation_2018} nomme l'\emph{analytical intelligence}~: l'intelligence qui organise le calcul, distincte de celle qui l'exécute et de celle qui conçoit les formules. Elle occupait le deuxième étage de la pyramide de Prony~--- les sept ou huit algébristes qui traduisaient les mathématiques en procédures élémentaires~--- mais la mécanisation l'avait considérablement alourdie, parce que les machines ne calculaient pas de la même manière que les humains, et que chaque machine calculait différemment des autres.

Les frictions étaient concrètes. Les frères Daniels, calculateurs anciens et irremplaçables pour la relecture d'épreuves, étaient jugés \og \emph{temperamentally unfitted to supervise subordinate staff}\fg{} et \og \emph{too stereotyped in their habits to adapt themselves to the use of machines}\fg{}. Miss~Stocks et Miss~Burroughs, affectées à la transformation des coordonnées héliocentriques en coordonnées géocentriques sur machines Brunsviga, nécessitèrent chacune trois mois de formation individuelle dispensée par le surintendant lui-même. Le nouveau système était, à terme, environ vingt pour cent moins cher que l'ancien~--- mais seulement parce que les opératrices étaient payées moins que les anciens calculateurs, pas parce que les machines avaient réduit le travail.

\paragraph{Le PLM~: \og d'abord organiser, ensuite mécaniser\fg{}}\mbox{}\\
Le même schéma se reproduisait dans un contexte radicalement différent. En 1929, Georges Bolle, polytechnicien, chef de la comptabilité de la Compagnie des chemins de fer Paris-Lyon-Méditerranée, publiait dans la \emph{Revue générale des chemins de fer} un article détaillé sur l'introduction des machines Hollerith pour le suivi du fret \parencite{bolle_utilisation_1929}. Bolle ne cachait pas son enthousiasme~: les avantages économiques des nouvelles machines étaient, à ses yeux, \og vraiment incalculables\fg{}. Mais l'essentiel de l'article ne portait pas sur les machines~; il portait sur le travail nécessaire \emph{avant} que les machines ne puissent fonctionner. Chaque détail du flux de travail devait être méticuleusement pensé à l'avance~: comment maximiser l'usage des quarante-cinq colonnes d'une carte Hollerith, comment numéroter les catégories de fret pour accélérer le codage, quels pictogrammes créer pour les marchandises les plus fréquemment transportées. \og \emph{Chaque détail doit être minutieusement examiné, discuté, pesé dans un travail de cette nature}\fg{}, écrivait Bolle, avant de formuler la devise qui résumait l'expérience du PLM~: \og \emph{D'abord organiser, ensuite mécaniser}\fg{}\footnote{Cité également par \textcite{couffignal_machines_1933}, p.~79, qui en fait un principe général de la mécanisation du calcul.}.

Comme au \emph{Nautical Almanac}, l'introduction des machines avait entraîné la centralisation du travail (à Paris), le recrutement d'une main-d'œuvre féminine bon marché, et un effort de supervision sans proportion avec celui de l'ancien système. Les opératrices perforaient trois cents cartes de quarante-cinq colonnes par jour, pendant six heures et demie au maximum, et pas plus de quatorze jours consécutifs par mois \parencite{bolle_utilisation_1929}. Une étude psychophysique conduite peu après par \textcite{lahy_selection_1931} sur les opératrices de machines Elliot-Fischer dans les compagnies ferroviaires françaises confirma que chaque calcul impliquait seize étapes distinctes~--- de l'insertion du papier dans la machine au réglage des compteurs en passant par la vérification des noms et des montants~--- et que l'attention requise était \og \emph{continue et concentrée}\fg{}, impossible à soutenir longtemps \og \emph{sans repos}\fg{}. L'effort que les machines avaient promis de supprimer~--- la fatigue mentale du calcul répétitif~--- n'avait pas disparu~; il s'était déplacé des calculateurs vers les opératrices (qui devaient maintenir une vigilance constante sur la machine) et vers les superviseurs (qui devaient concevoir les séquences d'opérations). \dimmark{D4}{*}%

\paragraph{Le résultat~: la mécanisation réorganise, elle ne remplace pas}\mbox{}\\
Les deux cas~--- l'un scientifique et britannique, l'autre industriel et français~--- convergent sur un résultat que ni Comrie ni Bolle n'avaient anticipé et que la littérature ultérieure n'a pas suffisamment relevé. La mécanisation du calcul ne réduit pas le nombre de travailleurs~: elle le maintient ou l'augmente, et elle en transforme le profil. Des calculateurs qualifiés, autonomes, travaillant à domicile, sont remplacés par des opératrices moins qualifiées, moins payées, supervisées de près, auxquelles il faut ajouter un surintendant dont le temps est désormais absorbé, à hauteur de vingt à trente pour cent, par la conception des procédures. Le coût total peut baisser~--- il baisse de vingt pour cent au \emph{Nautical Almanac}~--- mais cette baisse tient à la compression salariale, pas à l'économie de travail.

\dimmark{D4}{*}%
La devise de Bolle~--- \og d'abord organiser, ensuite mécaniser\fg{}~--- fait écho, sans que Bolle pût le savoir, aux six formes informationnelles de McCallum sur l'\emph{Erie Railroad}. Dans les deux cas, la technique matérielle (le télégraphe, la tabulatrice) n'est efficace qu'insérée dans un dispositif organisationnel qui la précède et la conditionne. Dans les deux cas, le coût de ce dispositif est sous-estimé par les promoteurs de la technique. Et dans les deux cas, ce coût est essentiellement un coût de travail qualifié~--- l'\emph{analytical intelligence} de Daston, la \og supervision\fg{} de Comrie, l'\og effort cérébral très laborieux\fg{} de Bolle~--- qui ne disparaît pas avec la machine mais se concentre sur un nombre restreint de personnes dont la charge cognitive augmente.

\emergence{D4~--- La mécanisation du calcul ne remplace pas le travail humain~: elle le réorganise. Le nombre de travailleurs se maintient ou augmente~; leur profil change (calculateurs qualifiés $\rightarrow$ opératrices moins payées + superviseurs plus chargés). Le coût de supervision (20--30\,\% du travail total) est une charge nouvelle créée par la mécanisation elle-même. \questionch{2}{D4~--- Cette restructuration sans suppression se conserve-t-elle avec le computing électronique~? L'ordinateur remplace-t-il enfin les calculateurs, ou réorganise-t-il le travail une deuxième fois~?}}

\emergence{D7~--- La devise de Bolle (\og d'abord organiser, ensuite mécaniser\fg{}) confirme, dans un contexte entièrement différent, le résultat du cas ferroviaire~: la technique informationnelle n'est efficace qu'insérée dans un dispositif organisationnel qui la précède. Le coût de ce dispositif (20--30\,\% du travail total au \emph{Nautical Almanac}) est un coût de coordination au sens de \textcite{coase_nature_1937}~--- mais un coût que la technique informationnelle \emph{crée} au lieu de le réduire. \questionch{2}{D7~--- La mécanisation du calcul augmente-t-elle les coûts de coordination à l'intérieur de la firme avant de les réduire~? Ce schéma se reproduit-il avec l'informatisation des années~1960--1980~?}}


%%% --- BILAN §1.3--§1.4 ---

\subsection*{Bilan~: deux mécanisations, deux logiques}\label{subsec:bilan_compute}

Les deux sections qui précèdent ont traversé, dans l'ordre chronologique, deux phénomènes que le mot \og mécanisation du calcul\fg{} tend à confondre. La calculatrice mécanique (§\,\ref{sec:calculating}) est un outil individuel qui accélère un geste élémentaire~--- l'addition, la multiplication~--- sans modifier le processus dans lequel ce geste s'insère, sans créer de mémoire externe, sans lier l'opérateur à un fournisseur particulier. Le système mécanographique (§\,\ref{sec:mechanography}) est une architecture intégrée~--- perforatrices, trieuses, tabulatrices, reproductrices~--- qui couple pour la première fois, dans un système opérationnel, le stockage de données sur un support lisible par la machine et le traitement en volume de ces données. La différence n'est pas de degré~; elle est de nature, et elle produit des résultats économiques opposés~: un marché concurrentiel dans un cas, un monopole par verrouillage dans l'autre~; un instrument invisible à l'extérieur de l'organisation dans un cas, un système qui restructure les organisations et leurs rapports avec leurs fournisseurs dans l'autre.

L'un et l'autre partagent cependant un trait commun, que le zoom sur le \emph{Nautical Almanac} et le PLM a fait apparaître~: la mécanisation du calcul ne supprime pas le travail humain~; elle le réorganise. Le nombre de travailleurs se maintient ou augmente~; leur profil change~--- des calculateurs qualifiés et autonomes sont remplacés par des opératrices moins payées et plus étroitement supervisées~; et le coût de supervision~--- ce que \textcite{daston_calculation_2018} nomme l'\emph{analytical intelligence}~--- croît au lieu de décroître. La mécanisation améliore la productivité du calcul, mais elle le fait en restructurant le travail, pas en le supprimant~--- un résultat que \textcite{reinstaller_market_2001}, \textcite{davies_womans_1982} et \textcite{strom_beyond_1992} documentent dans trois registres différents (quantitatif, genre, organisation). \textcite{yates_control_1989} montre que cette restructuration ne s'explique pas par la machine seule~: les systèmes de classement, les formulaires standardisés et les technologies de duplication interne précèdent les machines de bureau et conditionnent leur adoption.

\textcite{kooij_information_2022} est, dans la littérature sur les technologies à usage général, l'auteur qui accorde le plus d'attention à la trajectoire du calcul mécanique. Il mentionne l'Arithmomètre, les tabulatrices Hollerith, Babbage, et parle d'une \og \emph{First Mechanical Calculator Revolution}\fg{}. Mais il condense cent trente ans en un seul paragraphe de contexte, en route vers le transistor~; son cadre analytique complet (phases de vie de la technologie, cycles d'affaires, identification des brevets) n'est pas appliqué à cette période. Hollerith reçoit une demi-phrase~; Morse, Bell et Marconi reçoivent chacun un volume entier. Ce n'est pas un oubli~: c'est le signe que la mécanisation du calcul était trop petite, relativement à l'économie, pour que le cadre des technologies à usage général la capte comme un fait de premier ordre. L'industrie entière des machines de bureau pesait, à son pic de~1929, deux dixièmes de pour cent du produit national brut~; IBM seul en représentait un centième. À cette première relativisation, il faut en superposer une seconde~: l'industrie des machines à cartes perforées elle-même, dont IBM contrôlait~85,7\,\%, ne représentait jusqu'aux années~1930 qu'une part minime de la mécanographie totale. Les tâches comptables et bancaires courantes étaient assurées principalement par les machines comptables traditionnelles~--- Burroughs, Comptometer, Underwood Bookkeeping~---, pas par les tabulatrices Hollerith \parencite[p.~158]{cortada_before_1993}. \textcite{heide_punched_2009} confirme la lecture pour les marchés européens~: le système hollerithien ne s'imposa pas comme architecture universelle de la mécanographie de bureau, mais comme spécialisation d'un sous-segment~--- celui du dépouillement statistique en grand volume et de la gestion des flux dans les administrations et les grandes entreprises. Le monopole d'IBM portait donc, jusqu'à la veille de la Seconde Guerre mondiale, sur un sous-segment d'une industrie elle-même marginale dans l'économie américaine. L'énergie consommée par le calcul mécanique était, comme on l'a vu, 1,5\,\% du coût total d'une opération de mécanisation~--- un rapport qui rendait la question énergétique littéralement invisible aux acteurs comme aux observateurs\footnote{L'invisibilité énergétique ne signifie pas l'absence d'empreinte matérielle. \textcite{little_toxic_2014} a documenté que le site d'Endicott (New York), ville-usine d'IBM depuis les années~1920, est aujourd'hui contaminé par un panache de trichloroéthylène sur 300~acres, avec plus de 200~millions de dollars dépensés en décontamination. La continuité du site~--- de la mécanographie des années~1930 à la fabrication de composants électroniques de l'après-guerre~--- montre que la dimension matérielle traverse la trajectoire d'IBM bien au-delà du carton perforé. Voir également \textcite{lecuyer_clean_2017} sur les pollutions de la Silicon Valley et \textcite{chiu_dark_2011} sur celles de Taïwan.}.

Ce n'est qu'au moment où le calcul changera de substrat~--- du mécanique à l'électronique~--- et où son volume croîtra de plusieurs ordres de grandeur que les ratios commenceront à se modifier. Le tableau~\ref{tab:compute-dimensions} récapitule ce que la traversée des deux cas~--- la calculatrice mécanique et le système mécanographique~--- a permis d'établir pour chacune des dix dimensions de la grille de lecture.

\begin{table}[p]
\caption{La mécanisation du calcul au prisme des dix dimensions~: observations, incertitudes et questions ouvertes}\label{tab:compute-dimensions}
\centering\footnotesize
\renewcommand{\arraystretch}{1.15}
\begin{tabularx}{\textwidth}{@{}cXcp{4.5cm}@{}}
\toprule
\textbf{Dim.} & \textbf{Observations des cas \u00a7\,\ref{sec:calculating}--\ref{sec:mechanography}} & & \textbf{Question pour la suite} \\
\midrule
D1 & \emph{Calculatrice}~: traite des données internes, pas de support persistant. \emph{Mécanographie}~: la carte perforée est le premier support de données lisible par la machine, réutilisable et stockable. Le carton ne consomme pas d'énergie pour conserver l'information. & + & Le passage au support électronique modifie-t-il cette propriété~? \\[3pt]
D2 & \emph{Calculatrice}~: marché fragmenté ($\sim$100~firmes US, 1904). \emph{Mécanographie}~: monopole IBM à~85,7\,\% (DoJ, 1935) par le triptyque location + cartes exclusives + brevets (Peirce~1922, Tauschek~1930). Différentiel institutionnel~: licence à Powers~: 20\,\% US vs 5\,\% Allemagne \parencite{heide_punched_2009}. Quatre cycles de clôture identifiés par \textcite{mounier-kuhn_punched_2015}. & * & Ce mécanisme de verrouillage par le consommable se retrouve-t-il dans d'autres secteurs~? \\[3pt]
D3 & \emph{Calculatrice}~: bien d'équipement vendu ($\sim$\$220). \emph{Mécanographie}~: location + flux de consommables. Cartes = 33\,\% des revenus IBM (1929), 39\,\% pendant la dépression \parencite{cortada_before_1993}. & + & Le modèle location + consommable constitue-t-il un type de capital informationnel distinct~? \\[3pt]
D4 & La mécanisation ne semble pas réduire l'emploi de bureau~: il passe de 1--2\,\% (1870) à 10\,\% (1940) de la population active \parencite{cortada_before_1993}. Corrélation K/L~: 0,921 \parencite{reinstaller_market_2001}. Au \emph{Nautical Almanac}~: supervision = 20--30\,\% du travail \parencite{daston_calculation_2018}. Féminisation et compression salariale documentées par \textcite{davies_womans_1982} et \textcite{strom_beyond_1992}. & + & Le computing électronique modifie-t-il ce rapport~? \\[3pt]
D5 & Coût par opération en baisse de $\sim$3\,\%/an, 1850--1945 \parencite{nordhaus_two_2007}. Seuil client~: 1\,500~opérations/jour \parencite{cortada_before_1993}. \emph{Pas d'analyse d'élasticité disponible.} & ? & La baisse du coût engendre-t-elle un rebond de la demande de calcul~? \\[3pt]
D6 & Pas d'effet documenté sur les prix ou l'intégration des marchés. Le calcul agit à l'intérieur des organisations~; ses effets restent dans les budgets internes, qui ne sont généralement pas conservés. & ? & L'absence d'effet observable reflète-t-elle une différence de nature, ou un défaut de sources~? \\[3pt]
D7 & \og D'abord organiser, ensuite mécaniser\fg{} \parencite{bolle_utilisation_1929}. 16~étapes par calcul \parencite{lahy_selection_1931}. \textcite{yates_control_1989}~: les systèmes de classement précèdent les machines et conditionnent leur adoption. & + & La mécanisation augmente-t-elle les coûts de coordination avant de les réduire~? \\[3pt]
D8 & IBM revenus~: \$2,5M (1919) $\rightarrow$ \$45,3M (1940). Dividendes payés sans interruption pendant la dépression. & ? & Le financement du calcul suit-il un cycle propre~? \\[3pt]
D9 & Pas de chaîne d'extraction identifiée pour les machines de bureau. La carte est du carton~; la machine est de la mécanique de précision. & -- & La matérialité apparaît-elle avec le changement de substrat~? \\[3pt]
D10 & Powers disposait d'une technologie jugée supérieure en~1912--14 mais de capacités de fabrication et de commercialisation inférieures \parencite{cortada_before_1993, heide_punched_2009}. Carte 80~colonnes~: perforations rectangulaires délibérément brevetables \parencite{heide_punched_2009}. & + & L'avantage du premier entrant tient-il à la technique ou au modèle d'affaires~? \\[3pt]
\midrule
0-TER & £9 sur £607 \parencite{daston_calculation_2018}~: 1,5\,\% du coût total. L'énergie est un poste négligeable~; la question énergétique ne se pose pas à cette échelle. Industrie entière = 0,2\,\% du PNB (1929). & + & À quel moment et par quel mécanisme le coût énergétique du calcul change-t-il d'ordre de grandeur~? \\[3pt]
\bottomrule
\end{tabularx}
\renewcommand{\arraystretch}{1.0}
\end{table}


% =============================================================================